%========================================================================
\section{Conclusion \& projet professionnel}
% ⏱ ~3 min  |  Refermer le fil rouge, dire la trajectoire, finir net.
%========================================================================

\begin{frame}{Un fil rouge tenu sur trois missions}
  % 🎤 ~1 min 30. Boucler : le même geste d'ingénieur, décliné 3 fois.
  \begin{columns}[T]
    \column{.33\textwidth}\textbf{QA}\\ \footnotesize doc + pipeline $\to$ TNR reproductibles, savoir transmis
    \column{.33\textwidth}\textbf{Product}\\ \footnotesize critères implicites $\to$ tickets explicites, structure
    \column{.33\textwidth}\textbf{GTT}\\ \footnotesize docs épars $\to$ benchmark structuré $\to$ décision
  \end{columns}
  \vskip 10pt
  \begin{block}{La constante}
    \alert{Chaos $\to$ artefact reproductible et transmissible}, avec l'IA comme levier
    et le jugement d'ingénieur comme garde-fou.
  \end{block}
\end{frame}

\begin{frame}{Ce que ce stage a clarifié}
  % 🎤 ~1 min 30. Bilan technique / humain / pro. Trajectoire QA -> Product. CDI.
  \begin{itemize}
    \item \textbf{Technique} --- concevoir et mettre en prod une méthodo de test IA sur un site réel
    \item \textbf{Humain} --- défendre des choix devant des décideurs ; doser confiance et humilité
    \item \textbf{Professionnel} --- trajectoire vers le \alert{Produit} / la gestion de projet (CDI signé)
  \end{itemize}
  \vskip 8pt
  % 🎤 Mention bonne pratique INSA : un LLM a appuyé reformulation/mise en forme/relecture du rapport ;
  %    le fond --- analyses, arbitrages, conclusions --- reste personnel.
  \centerline{\footnotesize\textcolor{gray}{Continuer à intégrer l'IA comme un levier au service du métier.}}
\end{frame}
